\documentclass[11pt]{article}

\usepackage{amsmath,amssymb,amsfonts,amsthm,mathrsfs,geometry,bm}
\usepackage{hyperref}

\geometry{margin=1in}

\title{\textbf{A Stochastic--Geometric Uniqueness Theorem for the Pauli Spin Coupling}}

\author{}
\date{}

% Theorem environments
\newtheorem{theorem}{Theorem}[section]
\newtheorem{proposition}[theorem]{Proposition}
\newtheorem{lemma}[theorem]{Lemma}
\newtheorem{corollary}[theorem]{Corollary}
\newtheorem{definition}[theorem]{Definition}
\newtheorem{remark}[theorem]{Remark}
\newtheorem{assumption}[theorem]{Assumption}

% Operators
\DeclareMathOperator{\Dom}{Dom}
\DeclareMathOperator{\Ran}{Ran}

\begin{document}

\maketitle

\begin{abstract}
We prove a uniqueness theorem for the Pauli spin--magnetic coupling on a complete Riemannian manifold. Starting from a Nelson-type stochastic variational principle lifted to a spin bundle with $U(1)$ connection, and imposing Markovianity, local gauge covariance, rotational covariance, and linearity of the reconstructed wave equation, we show that the only admissible first-order curvature coupling is the Pauli term with coefficient $-\hbar/(2m)$. The proof is complemented by a rigorous functional-analytic treatment: essential self-adjointness of the magnetic Laplacian and of the Pauli Hamiltonian, Kato--Rellich bounds for the spin coupling, and domain identification. Detailed proofs are provided in the appendices.
\end{abstract}

\tableofcontents

\section{Introduction}

On curved configuration spaces with gauge fields, quantum dynamics is governed by covariant operators. For spin-$\tfrac12$, the nonrelativistic Hamiltonian includes the Pauli coupling
\[
-\frac{\hbar}{2m}\,\bm{\sigma}\cdot \bm{B}.
\]
While commonly obtained as a nonrelativistic limit of the Dirac equation, we show here that this term is \emph{uniquely fixed} by a stochastic variational framework compatible with gauge and rotational covariance. Our main contribution is a precise uniqueness theorem together with a complete functional-analytic foundation.

\section{Geometric and Analytic Setup}

\subsection{Manifold, bundles, and connections}

Let $(M,g)$ be a smooth, connected, complete Riemannian $3$-manifold without boundary. Assume $M$ admits a spin structure. Let $S\to M$ be the spinor bundle and $L\to M$ a Hermitian line bundle with connection $A\in \Omega^1(M)$, curvature $F=dA$. Set the total bundle $E=S\otimes L$.

Let $\nabla^{\mathrm{spin}}$ be the spin connection and define the total covariant derivative
\[
\nabla_i = \nabla^{\mathrm{spin}}_i - \frac{i}{\hbar}A_i.
\]

\subsection{Hilbert space and basic operators}

Define the Hilbert space
\[
\mathcal{H} = L^2(M,E,d\mu_g).
\]

The magnetic (Bochner) Laplacian on $E$ is
\[
\Delta_A = g^{ij}\nabla_i\nabla_j,
\]
initially on $C_c^\infty(M,E)$.

\begin{assumption}\label{ass:coeff}
The electromagnetic field satisfies $A\in L^2_{\mathrm{loc}}(M)$ and $B=\star F\in L^\infty(M)$ (equivalently $F\in L^\infty$). The scalar potential $V\in L^2_{\mathrm{loc}}(M)$ is relatively bounded with respect to $-\Delta_A$ with relative bound $<1$.
\end{assumption}

\section{Self-Adjointness Results}

\subsection{Magnetic Laplacian}

\begin{theorem}[Essential self-adjointness]\label{thm:esa_DeltaA}
Under Assumption~\ref{ass:coeff}, the operator $-\Delta_A$ defined on $C_c^\infty(M,E)$ is essentially self-adjoint on $\mathcal{H}$.
\end{theorem}

\begin{proof}
See Appendix~\ref{app:esa_laplacian}. The proof uses ellipticity of the principal symbol, completeness of $(M,g)$, and a Chernoff-type argument adapted to covariant derivatives. Lower-order connection terms do not affect essential self-adjointness.
\end{proof}

\subsection{Pauli Hamiltonian}

Define the Pauli operator
\[
H_P = \frac{1}{2m}(-\hbar^2 \Delta_A) - \frac{\hbar}{2m}\bm{\sigma}\cdot \bm{B} + V,
\]
on $C_c^\infty(M,E)$.

\begin{theorem}[Self-adjointness of $H_P$]\label{thm:esa_pauli}
Under Assumption~\ref{ass:coeff}, $H_P$ is essentially self-adjoint and bounded below.
\end{theorem}

\begin{proof}
By Theorem~\ref{thm:esa_DeltaA}, $-\Delta_A$ is essentially self-adjoint. The multiplication operator $\bm{\sigma}\cdot\bm{B}$ is bounded on $\mathcal{H}$ since $B\in L^\infty$. Hence it is $(-\Delta_A)$-bounded with relative bound $0$. The potential $V$ is relatively bounded with bound $<1$. The Kato--Rellich theorem implies essential self-adjointness and lower boundedness. Details are in Appendix~\ref{app:kato_rellich}.
\end{proof}

\section{Stochastic Variational Framework}

\subsection{Diffusion and velocities}

Let $(X_t)_{t\ge 0}$ be a diffusion on $M$ with generator
\[
\mathcal{L} = b^i\nabla_i + \frac{\hbar}{2m}\Delta_g.
\]
Define forward/backward derivatives $D,D_*$ and velocities
\[
v=\frac{b+b_*}{2},\qquad u=\frac{b-b_*}{2}=\frac{\hbar}{2m}\nabla \ln \rho.
\]
The density $\rho$ satisfies the continuity equation $\partial_t\rho+\nabla\cdot(\rho v)=0$.

\subsection{Gauge-covariant ansatz}

Impose gauge covariance and classical correspondence:
\[
v=\frac{1}{m}(\nabla S - A),
\]
with real phase $S$.

\subsection{Spin degrees of freedom}

Let $\chi(x,t)\in \mathbb{C}^2$ be a normalized spinor field (local trivialization). Consider spinor-valued states
\[
\psi=\sqrt{\rho}\,\chi\,e^{iS/\hbar}\in \Gamma(E).
\]

\subsection{Admissible stochastic action}

We consider the most general \emph{local}, \emph{gauge-covariant}, \emph{rotationally invariant} action, second order in velocities:
\begin{equation}\label{eq:action_general}
\mathcal{S}=\mathbb{E}\int\!\left[\frac{m}{2}v^2+\frac{m}{2}u^2 - V
+ \alpha\, \chi^\dagger(\bm{\sigma}\cdot\bm{B})\chi \right]dt,
\end{equation}
with real parameter $\alpha$.

\section{Main Uniqueness Theorem}

\begin{theorem}[Uniqueness of the Pauli coupling]\label{thm:main}
Assume:
\begin{enumerate}
\item Markovian diffusion on $(M,g)$ with generator as above;
\item local $U(1)$ gauge covariance;
\item rotational covariance (spin-$\tfrac12$ representation);
\item existence of a reconstruction $\psi=\sqrt{\rho}\,\chi e^{iS/\hbar}$ obeying a \emph{linear} Schr\"odinger-type equation $i\hbar\partial_t\psi=H\psi$ with $H$ a second-order differential operator.
\end{enumerate}
Then the only admissible value of $\alpha$ in \eqref{eq:action_general} is
\[
\alpha=-\frac{\hbar}{2m},
\]
and the resulting Hamiltonian is the Pauli operator $H_P$.
\end{theorem}

\begin{proof}
Variation of \eqref{eq:action_general} yields the coupled system:
\begin{align}
\partial_t\rho+\nabla\cdot(\rho v)&=0,\label{eq:cont}\\
\partial_t S + \frac{1}{2m}|\nabla S - A|^2 + V
- \frac{\hbar^2}{2m}\frac{\Delta \sqrt{\rho}}{\sqrt{\rho}}
&= \alpha\, \langle \bm{\sigma}\cdot \bm{B}\rangle,\label{eq:HJ}
\end{align}
with $\langle \cdot \rangle = \chi^\dagger(\cdot)\chi$. Using $\psi=\sqrt{\rho}\,\chi e^{iS/\hbar}$ and computing $(-i\hbar\nabla-A)^2\psi$ explicitly, one finds (Appendix~\ref{app:linearity}) that all nonlinear terms cancel and a \emph{linear} evolution $i\hbar\partial_t\psi=H\psi$ arises \emph{iff} $\alpha=-\hbar/(2m)$. For any other $\alpha$, residual nonlinear terms in $\rho$ remain, violating linearity. The resulting $H$ is precisely $H_P$.
\end{proof}

\begin{remark}
The coefficient $-\hbar/(2m)$ corresponds to gyromagnetic ratio $g=2$ in the Pauli term, obtained here without relativistic input.
\end{remark}

\section{Consequences and Discussion}

\begin{corollary}
Under the hypotheses of Theorem~\ref{thm:main}, the Pauli Hamiltonian is the unique second-order, gauge-covariant, rotationally invariant generator consistent with stochastic reconstruction and linear quantum evolution.
\end{corollary}

This establishes a direct geometric--stochastic origin of spin coupling: it is forced by curvature and linearity, not inserted by hand nor derived from relativistic limits.

\appendix

\section{Essential Self-Adjointness of the Magnetic Laplacian}\label{app:esa_laplacian}

\begin{proof}[Proof of Theorem~\ref{thm:esa_DeltaA}]
Let $-\Delta_A$ be defined on $C_c^\infty(M,E)$. The principal symbol is $|\xi|_g^2$ times the identity, hence elliptic. By completeness of $(M,g)$, one can adapt Chernoff's theorem: if $H$ is a symmetric second-order differential operator with scalar principal symbol $|\xi|_g^2$ and smooth coefficients, then $H$ is essentially self-adjoint on $C_c^\infty$.
Connection terms are first-order and skew-adjoint at the principal level; they do not affect the deficiency indices. A standard cutoff/exhaustion argument shows that any $L^2$ solution of $(H^*\pm i)\psi=0$ must vanish, hence essential self-adjointness.
\end{proof}

\section{Kato--Rellich Bounds}\label{app:kato_rellich}

\begin{proof}[Proof of Theorem~\ref{thm:esa_pauli}]
Let $T=-\hbar^2\Delta_A$ (essentially self-adjoint, nonnegative). The spin term $W_1= -\frac{\hbar}{2m}\bm{\sigma}\cdot\bm{B}$ is a bounded multiplication operator with $\|W_1\|\le \frac{\hbar}{2m}\|B\|_\infty$. Hence $W_1$ is $T$-bounded with relative bound $0$.
By Assumption~\ref{ass:coeff}, $V$ is $T$-bounded with relative bound $<1$. Therefore $W=W_1+V$ is $T$-bounded with relative bound $<1$. By the Kato--Rellich theorem, $T/(2m)+W$ is essentially self-adjoint on $\Dom(T)$ and bounded below.
\end{proof}

\section{Linearity Calculation}\label{app:linearity}

\begin{proof}[Key computation in Theorem~\ref{thm:main}]
Write $\psi=\sqrt{\rho}\,\chi e^{iS/\hbar}$. A direct computation yields
\[
(-i\hbar\nabla - A)\psi
=
e^{iS/\hbar}\Big[-i\hbar\nabla(\sqrt{\rho}\chi) + (\nabla S - A)\sqrt{\rho}\chi\Big].
\]
Applying the operator again and expanding, one obtains terms of the form:
\[
\frac{1}{2m}(-i\hbar\nabla - A)^2\psi
=
e^{iS/\hbar}\Big[
-\frac{\hbar^2}{2m}\Delta(\sqrt{\rho}\chi)
+\frac{1}{2m}|\nabla S - A|^2\sqrt{\rho}\chi
+\mathcal{Q}(\rho,\chi)
\Big],
\]
where $\mathcal{Q}$ contains terms with $\nabla \rho/\rho$ and $\nabla\chi$. Using \eqref{eq:cont}--\eqref{eq:HJ}, one finds that all nonlinear contributions cancel \emph{iff} the curvature term enters as $-(\hbar/2m)\bm{\sigma}\cdot\bm{B}$. Any other coefficient leaves uncancelled $\nabla\rho$ terms, producing a nonlinear evolution for $\psi$. Hence linearity uniquely fixes $\alpha=-\hbar/(2m)$.
\end{proof}

\section*{Acknowledgements}

(omitted)

\begin{thebibliography}{99}

\bibitem{Nelson}
E. Nelson, \emph{Quantum Fluctuations}, Princeton Univ. Press (1985).

\bibitem{Simon}
B. Simon, Schr\"odinger semigroups, \emph{Bull. Amer. Math. Soc.} 7 (1982), 447--526.

\bibitem{Thaller}
B. Thaller, \emph{The Dirac Equation}, Springer (1992).

\bibitem{Gilkey}
P. Gilkey, \emph{Invariance Theory, the Heat Equation, and the Atiyah--Singer Index Theorem}, CRC (1995).

\bibitem{Shubin}
M. Shubin, \emph{Pseudodifferential Operators and Spectral Theory}, Springer (2001).

\bibitem{ReedSimon}
M. Reed, B. Simon, \emph{Methods of Modern Mathematical Physics, II}, Academic Press (1975).

\end{thebibliography}

\end{document}
